\documentclass[11pt]{article}

% ============================================================
% Packages
% ============================================================
\usepackage[utf8]{inputenc}
\usepackage[margin=1in]{geometry}
\usepackage{amsmath,amssymb,amsthm}
\usepackage{graphicx}
\usepackage{booktabs}
\usepackage{xcolor}
\usepackage{hyperref}
\usepackage{tikz}
\usepackage{enumitem}
\usepackage{pgfplots}

\pgfplotsset{compat=1.18}

% ============================================================
% Document Formatting
% ============================================================
\setlength{\parskip}{0.6em}
\setlength{\parindent}{0pt}

% ============================================================
% List Formatting
% ============================================================
\setlist[enumerate]{
    leftmargin=2em,
    itemsep=0.25em
}

% ============================================================
% Solution Environment
% ============================================================
\newenvironment{solution}
    {\par\noindent\textbf{Solution.}\par}
    {\par}

% ============================================================
% Link Styling
% ============================================================
\hypersetup{
    colorlinks=true,
    linkcolor=blue,
    urlcolor=cyan
}

% ============================================================
% Metadata
% ============================================================
\title{Ordinary Differential Equations}
\author{Chris Ritter}
\date{October 20, 2026}

% ============================================================
% Document
% ============================================================


\begin{document}

\maketitle
\tableofcontents
\newpage

% ============================================================
\section{The Method of Variation of Parameters}
% ============================================================

Consider the nonhomogeneous second-order linear equation
\[
L[y]:=y''+p(t)y'+q(t)y=g(t).
\]
As before, its general solution has the form \(y=y_h+y_p\), where
\(L[y_h]=0\) and \(L[y_p]=g\). Undetermined coefficients is convenient
for polynomial, exponential, and trigonometric forcing terms, but
forcing terms such as \(\tan t\) or \(\csc t\) do not lead to a finite
trial of that kind. \textbf{Variation of parameters} constructs a
particular solution using a known fundamental pair \(y_1,y_2\) for
\(L[y]=0\).

The method works on any interval where \(p,q,g\) are continuous and
\(y_1,y_2\) are linearly independent. It is important first to divide
through by the coefficient of \(y''\) if the equation is not already
in the normalized form above.

\subsection{Derivation of the Formula}

The homogeneous solution is
\[
y_h=C_1y_1+C_2y_2.
\]
Replace the constants by functions and seek
\[
y_p=u_1(t)y_1(t)+u_2(t)y_2(t).
\]
Choose the auxiliary condition
\[
u_1'y_1+u_2'y_2=0.
\]
Then \(y_p'=u_1y_1'+u_2y_2'\), and differentiating once more gives
\[
y_p''=u_1'y_1'+u_2'y_2'+u_1y_1''+u_2y_2''.
\]
Because \(L[y_1]=L[y_2]=0\), substitution into \(L[y_p]=g\) leaves
\[
u_1'y_1'+u_2'y_2'=g.
\]
Thus the unknown derivatives satisfy the coefficient system
\[
\begin{pmatrix}y_1&y_2\\y_1'&y_2'\end{pmatrix}
\begin{pmatrix}u_1'\\u_2'\end{pmatrix}
=\begin{pmatrix}0\\g\end{pmatrix}.
\]
Its determinant is the \textbf{Wronskian}
\[
W=y_1y_2'-y_1'y_2.
\]
Since the solutions form a fundamental pair, \(W\ne0\) on the chosen
interval. Solving the system and integrating gives
\[
\boxed{u_1'=-\frac{y_2g}{W},\qquad
u_2'=\frac{y_1g}{W},\qquad
y_p=-y_1\int\frac{y_2g}{W}\,dt
+y_2\int\frac{y_1g}{W}\,dt.}
\]
Constants of integration only add homogeneous terms, so they can be
set to zero when finding \(y_p\).

\subsection{Procedure}

\begin{enumerate}
\item Normalize the equation so the coefficient of \(y''\) is one.
\item Solve the homogeneous equation and choose independent
      solutions \(y_1,y_2\).
\item Compute \(W=y_1y_2'-y_1'y_2\).
\item Integrate \(u_1'=-y_2g/W\) and \(u_2'=y_1g/W\).
\item Form \(y_p=u_1y_1+u_2y_2\), then add \(y_h\).
\end{enumerate}

\subsection{Example: Forcing by \texorpdfstring{$3\csc t$}{3 csc t}}

Find the general solution of
\[
y''+4y=3\csc t.
\]

\begin{solution}
Work on an interval \((k\pi,(k+1)\pi)\), where \(\sin t\ne0\).
The homogeneous characteristic equation \(r^2+4=0\) gives
\[
y_1=\cos2t,\qquad y_2=\sin2t,\qquad
W=2\cos^22t+2\sin^22t=2.
\]
Consequently,
\[
u_1'=-\frac{3\sin2t}{2\sin t}=-3\cos t,
\qquad
u_1=-3\sin t,
\]
and, using \(\cos2t=1-2\sin^2t\),
\[
\begin{aligned}
u_2'
&=\frac{3\cos2t}{2\sin t}
=\frac32\csc t-3\sin t,\\
u_2
&=\frac32\ln\left|\tan\frac{t}{2}\right|+3\cos t.
\end{aligned}
\]
The particular solution simplifies via
\(\sin2t\cos t-\cos2t\sin t=\sin t\):
\[
\begin{aligned}
y_p
&=-3\sin t\cos2t
+\left(\frac32\ln\left|\tan\frac{t}{2}\right|
  +3\cos t\right)\sin2t\\
&=3\sin t+\frac32\sin2t
  \ln\left|\tan\frac{t}{2}\right|.
\end{aligned}
\]
Therefore
\[
\boxed{y=C_1\cos2t+C_2\sin2t+3\sin t
+\frac32\sin2t\ln\left|\tan\frac{t}{2}\right|.}
\]
The formula is interpreted separately on each interval that does not
contain a zero of \(\sin t\).
\end{solution}

\subsection{Example: Forcing by \texorpdfstring{$\tan t$}{tan t}}

Find the general solution of
\[
y''+y=\tan t.
\]

\begin{solution}
Choose an interval where \(\cos t\ne0\), for example
\((-\pi/2,\pi/2)\). The homogeneous equation has fundamental
solutions
\[
y_1=\cos t,\qquad y_2=\sin t,\qquad W=1.
\]
The parameter derivatives are
\[
\begin{aligned}
u_1'&=-\sin t\tan t=\cos t-\sec t,\\
u_2'&=\cos t\tan t=\sin t.
\end{aligned}
\]
Thus we may take
\[
u_1=\sin t-\ln|\sec t+\tan t|,
\qquad
u_2=-\cos t.
\]
The mixed sine-cosine terms cancel in \(u_1y_1+u_2y_2\), leaving
\[
y_p=-\cos t\ln|\sec t+\tan t|.
\]
Hence
\[
\boxed{y=C_1\cos t+C_2\sin t
-\cos t\ln|\sec t+\tan t|.}
\]
The solution is valid on any interval that contains no pole of
\(\tan t\).
\end{solution}

\subsection{Key Points}

\begin{itemize}
\item The Wronskian must be nonzero on the interval of work.
\item The forcing term in the formula is the one from the normalized
      equation \(y''+py'+qy=g\).
\item Different antiderivatives for \(u_1,u_2\) produce particular
      solutions that differ only by homogeneous terms.
\item The method can require harder integrals than undetermined
      coefficients, but it handles a much wider range of forcing terms.
\end{itemize}

\end{document}
